%%% FIELD model-class-construction
The definition has the following extensionally equivalent two-line form. First,
\[
\begin{aligned}
\mathfrak M_{\mathrm{ocsi}}(G,\mathcal C)
:=\{\vartheta\in\Theta(G):{}&\text{ every latent vertex is a parentless root with exactly two observed children,}\\
&\text{the latent roots are mutually independent, every primitive CPT entry is strictly positive,}\\
&\text{every activation variable in }\mathcal C\text{ is observed, and every whole-row equality in }\mathcal C\text{ holds}\}.
\end{aligned}
\]
Second, let \((G,L_{C_0})\) be a maximal-regular labeled-DAG input whose controls \(C_0\subseteq V\) are observed roots, and let \(\mathcal C(L_{C_0})\) be the whole-row compilation that, whenever \(a\to i\in L_{C_0}(c)\), equates the complete conditional vectors of child \(i\) across values of parent \(a\) under the activation \(C_i=c|_{C_i}\), where \(C_i=\operatorname{pa}_G(i)\cap C_0\). Put \(\mathbf d=(|\mathcal X_j|:j\in V\cup U)\), write \(\mathfrak M_{\mathrm{MJEK}}(G,L_{C_0})\) for the published compatible labeled-DAG CPT class, and define
\[
\Theta^{\mathrm{sm},+}_{\mathbf d}(G)
:=\{\vartheta\in\Theta(G):\text{the fixed finite state spaces, latent-root, two-child-latent, independent-latent, and strict primitive-positivity restrictions hold}\}.
\]
On this compiled input, the same set is
\[
\mathfrak M_{\mathrm{ocsi}}\!\left(G,\mathcal C(L_{C_0})\right)
=\mathfrak M_{\mathrm{MJEK}}(G,L_{C_0})\cap\Theta^{\mathrm{sm},+}_{\mathbf d}(G).
\]
%%% FIELD mjek-criterion-statement
Let \(C_0\) be a set of observed root vertices of a finite labeled DAG \(G\), let \(X\) and \(Y\) be disjoint vertex sets outside \(C_0\), and let \(L_{C_0}\) be maximal-regular with respect to \(G\) in the sense of Mokhtarian--Jamshidi--Etesami--Kiyavash. In their positive-observed-law structural-equation-model class, \(P_{\mathbf x}(\mathbf y)\) is identifiable from \((G,L_{C_0})\) if and only if it is identifiable from every context DAG
\[
G_c=\bigl((V\setminus C_0)\cup U,\, (E(G)\cap(((V\setminus C_0)\cup U)\times((V\setminus C_0)\cup U)))\setminus L_{C_0}(c)\bigr),
\qquad c\in\mathcal X_{C_0}.
\]
When every context query is identified by an observed-law formula \(F_c\), their Algorithm 1 returns
\[
\sum_{c\in\mathcal X_{C_0}}P(C_0=c)F_c\!\left(P(V\setminus C_0\mid C_0=c)\right).
\]
%%% FIELD root-control-bridge-statement
Let \((G,L_{C_0},X,Y,\mathbf x,\mathbf y)\) be a maximal-regular observed-root input of Mokhtarian--Jamshidi--Etesami--Kiyavash, with the used hypotheses written explicitly as follows: \(C_0\subseteq V\) consists of observed roots; \(X,Y\subseteq V\setminus C_0\) are disjoint; \(L_{C_0}(c)\) is the label set at \(c\in\mathcal X_{C_0}\); and, for every labeled edge \(a\to i\), the set \(C_i=\operatorname{pa}_G(i)\cap C_0\) is nonempty and
\[
c|_{C_i}=c'|_{C_i}
\quad\Longrightarrow\quad
\bigl(a\to i\in L_{C_0}(c)\Longleftrightarrow a\to i\in L_{C_0}(c')\bigr).
\]
Compile a label \(a\to i\in L_{C_0}(c)\) into the whole-row CSI equating the complete conditional vectors of child \(i\) across values of parent \(a\), activated at \(C_i=c_i\), where \(c_i=c|_{C_i}\). Denote the compiled list by \(\mathcal C(L_{C_0})\). Put \(\mathbf d=(|\mathcal X_j|:j\in V\cup U)\). Let \(\Theta^{\mathrm{sm},+}_{\mathbf d}(G)\) be this note's fixed-finite-state, latent-root, two-child-latent, independent-latent, strictly primitive-positive CPT class before imposing CSI, and let \(\mathfrak M_{\mathrm{MJEK}}(G,L_{C_0})\) denote the published compatible labeled-DAG class in the same stochastic-CPT representation. Then
\[
\mathfrak M_{\mathrm{ocsi}}\!\left(G,\mathcal C(L_{C_0})\right)
=
\mathfrak M_{\mathrm{MJEK}}(G,L_{C_0})
\cap \Theta^{\mathrm{sm},+}_{\mathbf d}(G).
\]
For every \(c\), the complete active compiled trace \(t_c\) has
\[
H_{k(c),t_c}[(V\setminus C_0)\cup U]=G_c,
\]
where \(k(c)\) is the activation-partition cell containing context \(c\); and the compiled quotient \(\mathcal R_i\) is exactly the connected row-equality quotient generated by the labels of child \(i\). If every \(G_c\) identifies the context query with formula \(F_c\), the compiled SUCCESS derivations assemble the same formula as the published algorithm,
\[
E(P_{\mathrm{obs}})
=\sum_{c\in\mathcal X_{C_0}}P(C_0=c)
F_c\!\left(P(V\setminus C_0\mid C_0=c)\right)
=P_{\mathbf x}(\mathbf y),
\]
and every displayed conditioning divisor is positive. This transfers the published identification implication to the intersection above. It does not assert that failure in one \(G_c\) implies nonidentification on the narrower fixed-cardinality positive class. Finally, \(\mathsf E_5\) is admitted by this note's input syntax but not by the observed-root-control syntax: its activation variable \(X\) is observed and satisfies \(U\to X\), so the activation \(X=0\) is endogenous.
%%% FIELD root-control-bridge-proof
Write \(E(G)\) for the edge set of \(G\) and
\[
\mathcal X_A:=\prod_{j\in A}\mathcal X_j
\]
for the finite state space of a vertex set \(A\). We first check that the compilation is well defined. For a child \(i\), put
\[
C_i:=\operatorname{pa}_G(i)\cap C_0.
\]
The stated regularity hypothesis says that \(C_i\neq\varnothing\) whenever an edge into \(i\) is labeled. The stated maximality hypothesis says that if \(c,c'\in\mathcal X_{C_0}\) have \(c|_{C_i}=c'|_{C_i}\), then
\[
a\to i\in L_{C_0}(c)
\quad\Longleftrightarrow\quad
a\to i\in L_{C_0}(c').
\]
Consequently membership of \(a\to i\) in the active label set is a function only of \(c_i=c|_{C_i}\), so one may compile each distinct pair \((a\to i,c_i)\) exactly once.

More explicitly, for complete parent assignments \(s,s'\in\mathcal X_{\operatorname{pa}_G(i)}\), the compiled relation activated at \(C_i=c_i\) is
\[
s|_{\operatorname{pa}_G(i)\setminus\{a\}}
=s'|_{\operatorname{pa}_G(i)\setminus\{a\}},
\qquad
s|_{C_i}=s'|_{C_i}=c_i
\quad\Longrightarrow\quad
\theta_i(\,\cdot\mid s)=\theta_i(\,\cdot\mid s').
\tag{1}
\]
The label semantics say precisely that, in context \(c\), the complete primitive mechanism of \(i\) does not depend on parent \(a\). Such mechanism invariance implies (1) after taking the distribution of the mechanism's independent response noise. Conversely, order the finite states of \(i\), introduce an independent \(E_i\sim\operatorname{Unif}(0,1)\), and represent a CPT by
\[
f_i(s,E_i):=\min\left\{z:\sum_{z'\le z}\theta_i(z'\mid s)\ge E_i\right\}.
\]
If (1) holds, the cumulative sums, and hence \(f_i(s,e)\), are identical for all \(e\in(0,1)\) when only \(s_a\) changes in the active context. Thus the CPT has a label-compatible structural representation without changing any observational or interventional distribution. Therefore, in the stochastic-CPT representation, a table satisfies all compiled whole-row relations if and only if it is compatible with every label in \(L_{C_0}\). By \(\texttt{def:observed-csi-model-class}\), membership in \(\mathfrak M_{\mathrm{ocsi}}(G,\mathcal C(L_{C_0}))\) additionally imposes exactly the fixed finite state spaces, latent-root restriction, two-child-latent restriction, independent latent roots, and strict primitive positivity collected in \(\Theta^{\mathrm{sm},+}_{\mathbf d}(G)\). Therefore, in both directions,
\[
\theta\in\mathfrak M_{\mathrm{ocsi}}\!\left(G,\mathcal C(L_{C_0})\right)
\quad\Longleftrightarrow\quad
\theta\in\mathfrak M_{\mathrm{MJEK}}(G,L_{C_0})
\text{ and }
\theta\in\Theta^{\mathrm{sm},+}_{\mathbf d}(G),
\]
which proves the asserted class equality and verifies the second line of the replacement construction directly from the first. Notice that this is an intersection with the published class, not a replacement of that class by a new one.

We next compare the finite graphical syntax. For \(c\in\mathcal X_{C_0}\), let \(k(c)\) denote the fiber of the compiled joint activation truth-table map that contains assignments with control value \(c\). This is well defined even if two full control assignments lie in the same fiber: maximality then gives the same active labels. Let \(t_c\) enumerate once each edge having a compiled relation active on \(k(c)\). The preceding equivalence gives
\[
\{e:e\text{ occurs in }t_c\}=L_{C_0}(c).
\tag{2}
\]
Every entry is CSI-authorized on \(k(c)\), so \(t_c\) is an authorized trace. By \(\texttt{def:context-graph}\) and (2),
\[
H_{k(c),t_c}=G\setminus L_{C_0}(c).
\]
Removing the observed control roots therefore gives
\[
\begin{aligned}
H_{k(c),t_c}[(V\setminus C_0)\cup U]
&=\Bigl((V\setminus C_0)\cup U,
      \bigl(E(G)\setminus L_{C_0}(c)\bigr)
      \cap(((V\setminus C_0)\cup U)\times((V\setminus C_0)\cup U))\Bigr)\\
&=\Bigl((V\setminus C_0)\cup U,
      \bigl(E(G)\cap(((V\setminus C_0)\cup U)\times((V\setminus C_0)\cup U))\bigr)
      \setminus L_{C_0}(c)\Bigr)\\
&=G_c.
\end{aligned}
\tag{3}
\]
Thus the context graphs are identical, not merely Markov equivalent.

For the row quotient, define a finite undirected graph \(B_i^L\) whose vertices are the complete assignments in \(\mathcal X_{\operatorname{pa}_G(i)}\), with an edge between \(s\) and \(s'\) exactly when (1) is one of the equalities generated by a label. By construction of \(\mathcal C(L_{C_0})\), its primitive row-equality graph for child \(i\) is literally \(B_i^L\). By \(\texttt{def:row-quotient}\),
\[
\mathcal R_i=\pi_0(B_i^L),
\tag{4}
\]
the set of connected components. Equation (4) is the connected equality closure of the published labels, and \(\texttt{thm:canonical-row-quotient}\) ensures that passing to this quotient changes neither the compatible table class nor its observed or interventional maps.

It remains to compare returned formulas and to separate that comparison from a converse claim. Fix
\(
\theta\in\mathfrak M_{\mathrm{ocsi}}(G,\mathcal C(L_{C_0})).
\)
Strict primitive positivity and finiteness imply, for every complete observed assignment \(v\),
\[
P_\theta(V=v)
=\sum_{u\in\mathcal X_U}
  \prod_{j\in V\cup U}
  \theta_j\!\left((v,u)_j\mid(v,u)_{\operatorname{pa}_G(j)}\right)>0.
\tag{5}
\]
The sum is finite and nonempty, and each summand is a product of strictly positive coordinates. Summing (5) over assignments extending \(C_0=c\) gives
\[
P_\theta(C_0=c)>0.
\tag{6}
\]
Thus every conditional observed law \(P_\theta(V\setminus C_0\mid C_0=c)\) is defined by a divisor certified positive on the model class.

Because every member of \(C_0\) is a root and \(X\cap C_0=\varnothing\), truncated factorization gives
\[
P_{\theta,\mathbf x}(C_0=c)=P_\theta(C_0=c).
\tag{7}
\]
Conditioning that same factorization on \(C_0=c\) fixes the root mechanisms and deletes the label-active edges. In view of (3), the resulting postintervention law is the context-model law on \(G_c\), and hence
\[
P_{\theta,\mathbf x}(\mathbf y\mid C_0=c)
=P^{(c)}_{\theta,\mathbf x}(\mathbf y).
\tag{8}
\]
Suppose ordinary identification in every \(G_c\) returns \(F_c\). The cited Shpitser--Pearl output theorem and (8) give
\[
F_c\!\left(P_\theta(V\setminus C_0\mid C_0=c)\right)
=P^{(c)}_{\theta,\mathbf x}(\mathbf y)
=P_{\theta,\mathbf x}(\mathbf y\mid C_0=c).
\tag{9}
\]
Starting from observed cells, the conditioning in (9) uses the positive denominator (6); the remaining marginalization, conditioning, district restriction, and leaf-removal nodes of the ordinary derivation are constructors allowed by \(\texttt{def:success-derivation}\), and every edge deletion is justified by \(t_c\). Therefore these are compiled SUCCESS derivations. Total probability, followed by (7) and (9), yields
\[
\begin{aligned}
P_{\theta,\mathbf x}(\mathbf y)
&=\sum_{c\in\mathcal X_{C_0}}
  P_{\theta,\mathbf x}(\mathbf y\mid C_0=c)
  P_{\theta,\mathbf x}(C_0=c)\\
&=\sum_{c\in\mathcal X_{C_0}}
  P_\theta(C_0=c)
  F_c\!\left(P_\theta(V\setminus C_0\mid C_0=c)\right)\\
&=E(P_\theta(V)).
\end{aligned}
\tag{10}
\]
The sum is finite, and all divisors have recorded positivity judgments, so \(E\in\mathscr E(P_{\mathrm{obs}})\). This is exactly the assembly in the cited Mokhtarian--Jamshidi--Etesami--Kiyavash algorithm. Its validity on this note's class also follows from \(\texttt{thm:success-soundness}\).

The logical direction transferred in (10) is the correct one. The class equality gives
\[
\mathfrak M_{\mathrm{ocsi}}\!\left(G,\mathcal C(L_{C_0})\right)
\subseteq \mathfrak M_{\mathrm{MJEK}}(G,L_{C_0}).
\tag{11}
\]
An observed-law functional valid on the right side of (11) remains valid on the left side. By contrast, a pair witnessing nonidentification on the right need not have the fixed latent cardinalities or strictly positive primitive CPT coordinates required on the left. Thus the published necessity proof does not, merely by restriction, supply a countermodel pair in this note's class. No converse for the intersection is used or claimed.

Finally, in \(\mathsf E_5\) the activation predicate is \(X=0\), the activated whole-row relation is
\[
P(Z=0\mid X=0,T=0)=\frac12
=P(Z=0\mid X=0,T=1),
\]
with the complementary \(Z=1\) probabilities equal as well. Hence it is a whole-row CSI with an observed activation variable, as required by \(\texttt{ass:observed-activation}\) and \(\texttt{ass:whole-row-csi}\). But the displayed graph contains \(U\to X\), so \(X\) is not a root. Indeed, every observed vertex of \(\mathsf E_5\) has a parent: \(U\to X\), \(T\to Z\) and \(X\to Z\), and \(U,T,Z\to Y\). Therefore its nonconstant activation cannot be represented by a control set contained in the observed roots. This proves strict enlargement of the permitted input syntax beyond observed-root controls.
%%% FIELD oeq-narrowed-claim
Can \(\mathsf H_{\mathrm{sat}}\) be completed for every finite \((G,\mathcal C,\mathbf x,\mathbf y)\) so that exhaustive sound context-factor recovery either assembles a total \(E\in\mathscr E(P_{\mathrm{obs}})\), or produces retained partial-trace rooted C-forests whose jointly quotient-supported multi-block contraction has a kernel direction with zero full observed increment and nonzero query increment, with the binary-sink construction as its exact one-block base case? The resulting procedure must be sound and complete. On maximal-regular observed-root inputs, the required comparison is the compilation proved in \(\texttt{thm:root-control-class-bridge}\): identical connected row quotients and context graphs, together with agreement with the Mokhtarian--Jamshidi--Etesami--Kiyavash returned formula whenever their per-context identification criterion succeeds. A semantic nonidentification converse on this note's narrower fixed-cardinality, strictly primitive-positive class is required only if supported by a separate countermodel construction in that class. The permitted input syntax must still strictly include observed endogenous activations, as witnessed by \(\mathsf E_5\).
%%% FIELD oeq-replacement-claim
Can \(\mathsf H_{\mathrm{sat}}\) be completed for every finite \((G,\mathcal C,\mathbf x,\mathbf y)\) so that exhaustive sound context-factor recovery either assembles a total \(E\in\mathscr E(P_{\mathrm{obs}})\), or produces retained partial-trace rooted C-forests whose jointly quotient-supported multi-block contraction has a kernel direction with zero full observed increment and nonzero query increment, with the binary-sink construction as its exact one-block base case? The resulting procedure must be sound and complete. On maximal-regular observed-root inputs, the required comparison is the compilation proved in \(\texttt{thm:root-control-class-bridge}\): identical connected row quotients and context graphs, together with agreement with the Mokhtarian--Jamshidi--Etesami--Kiyavash returned formula whenever their per-context identification criterion succeeds. A semantic nonidentification converse on this note's narrower fixed-cardinality, strictly primitive-positive class is required only if supported by a separate countermodel construction in that class. The permitted input syntax must still strictly include observed endogenous activations, as witnessed by \(\mathsf E_5\).
%%% FIELD oeq-partial-result
The exact row quotient, sound SUCCESS assembly, fixed-nuisance binary-sink affine reduction, corrected row-space/nullspace alternative, and positive two-model certificate are proved. The forest-to-kernel lift yields the verifier's canonical forest-supported rational null vector. The five-node example proves that maximal active deletion can erase the obstruction, while the bow example proves that a local forest can be annihilated by a shared observed reference context. The root-control bridge proves exact compilation of maximal-regular observed-root inputs, equality with the published class intersected with this note's fixed finite strictly primitive-positive CPT class, identical context graphs and row quotient, and agreement of successful returned formulas; it deliberately transfers no nonidentification converse to that intersection. Thus the remaining theorem must retain partial traces and solve a jointly shared-row multi-block contraction problem whose separating direction changes the full query.
%%% FIELD oeq-obligation-partial-result
The exact row quotient, sound SUCCESS assembly, fixed-nuisance binary-sink affine reduction, corrected row-space/nullspace alternative, and positive two-model certificate are proved. The forest-to-kernel lift yields the verifier's canonical forest-supported rational null vector. The five-node example proves that maximal active deletion can erase the obstruction, while the bow example proves that a local forest can be annihilated by a shared observed reference context. The root-control bridge proves exact compilation of maximal-regular observed-root inputs, equality with the published class intersected with this note's fixed finite strictly primitive-positive CPT class, identical context graphs and row quotient, and agreement of successful returned formulas; it deliberately transfers no nonidentification converse to that intersection. Thus the remaining theorem must retain partial traces and solve a jointly shared-row multi-block contraction problem whose separating direction changes the full query.
